\section{Conclusion}\label{sec:conclusion}
Within this work, we introduce the idea of \textit{implicit assumptions} as a novel, unified perspective on the difficulties in reproducing, replicating, and reusing data analysis artifacts. 
Based on this perspective, we envision a set of techniques to address these difficulties with the help of static code analysis, and to enrich these artifacts with the necessary information to dynamically ensure their reproducibility and reusability. We combine this with an explorative implementation on top of \textit{flowR} which underlines the feasibility of our approach.
Yet, these categories are by no means exhaustive. For example, research artifacts may rely on undocumented, manual setup steps or interactivity from the user,
which hinder automated execution, and reduce the reproducibility of the
results.
Furthermore, such techniques are no universal solution as inferring these implicit assumptions is inherently limited by Rice's theorem~\cite{rice1953classes}.
Hence, we require thorough analyses of reproducible and irreproducible artifacts to identify common patterns to use as heuristics for these techniques. 
Additionally, we need to consider a human-in-the-loop approach if the detection of an implicit assumption works, but can not be automatically resolved.
However, even in these cases, we believe that the mere detection of implicit assumptions is still a valuable contribution, offering an interesting and valuable research direction.
\section{Future Plans}\label{sec:future-plans}
So far, we have already started working on a proof-of-concept implementation of our idea on top of \textit{flowR}~\cite{floopsla} and intend to complete this implementation by adding best-effort support for the detection of all implicit assumptions mentioned in \cref{sec:overview}.
To allow for additional constraints by users, we intend on adding a combination of a domain-specific language and the ability to detect violations of common assumptions as with \href{https://rdrr.io/r/base/stopifnot.html}{\texttt{stopifnot}} in~R. 
Subsequently, we plan to evaluate the usefulness of our approach in three ways:
\begin{itemize}[leftmargin=*,nosep]
   \item Take non-executable artifacts from public repositories and published studies~(e.g., those found by~\citeauthor{trisovic_largescale_2022}~\cite{trisovic_largescale_2022}) and check how many of them can be made executable with our approach, comparing this to the related work~(cf.~\cite{trisovic_largescale_2022,DBLP:conf/msr/IslamAW24,donat2025r4r}).
   \item Take executable artifacts and use fuzzing on their datasets to produce problematic and unproblematic variations and to get an understanding of the effectiveness of the inferred constraints to prevent problems in reusing the artifact with different data.
   \item Conduct a user study with R users to evaluate the usefulness of the inferred constraints as documentation for understanding and reusing the artifact by measuring their performance in understanding and reusing the artifact with and without the constraints.
\end{itemize}
